% dynamicalwh_figs.tex   2002-05-03
%%%%%%%%%%%%%%%%%%%%%%%%%%%%%%%%%%%%%%%%%%%%%%%%%%%%%%%%%%%%%%%%%%%%%%%%%%%%%%
%  FIGs
%%%%%%%%%%%%%%%%%%%%%%%%%%%%%%%%%%%%%%%%%%%%%%%%%%%%%%%%%%%%%%%%%%%%%%%%%%%%%%

%>>>>>>>>>>>>>>>>>>>>>>>>>>>>>>>>>>>>>>> Fig.\ref{fig1_grid}
\begin{figure}[ht]
\setlength{\unitlength}{1mm}
\begin{picture}(80,80)
\put(0,0){\epsfxsize=70mm \epsffile{fig1.eps} }
\end{picture}
\caption[quartic]{ Numerical grid structure.  Initial data are given on null
hypersurfaces $\Sigma_\pm$ ($x^\mp=0$, $x^\pm>0$) and their intersection $S$. }
\label{fig1_grid}
\end{figure}
%<<<<<<<<<<<<<<<<<<<<<<<<<<<<<<<<<<<<<<<     \ref{fig1_grid}


%>>>>>>>>>>>>>>>>>>>>>>>>>>>>>>>>>>>>>>> Fig.\ref{fig2_convergence}
\begin{figure}[t]
\setlength{\unitlength}{1mm}
\begin{picture}(85,50)
\put(0,0){\epsfxsize=80mm \epsffile{fig2.eps} }
\end{picture}
\caption[quartic]{ Convergence behaviour of the code for exact static wormhole
initial data. The location of the trapping horizon $\vartheta_-=0$ is plotted
for several resolutions labelled by the number of grid
points for $x^+=[0,20]$. We see that
numerical truncation error eventually destroys the static configuration. }
\label{fig2_convergence}
\end{figure}
%<<<<<<<<<<<<<<<<<<<<<<<<<<<<<<<<<<<<<<<     \ref{fig2_convergence}

%>>>>>>>>>>>>>>>>>>>>>>>>>>>>>>>>>>>>>>> Fig.\ref{fig3_config}
\begin{figure}[t]
\setlength{\unitlength}{1mm}
\begin{picture}(85,140)
\put(5,65){\epsfxsize=70mm \epsffile{fig3a.eps} }
\put(5,0){\epsfxsize=70mm \epsffile{fig3b.eps} }
\end{picture}
\caption[quartic]{ Static wormhole configuration obtained with the highest
resolution calculation: (a) expansion $\vartheta_+$ and (b) local gravitational
mass-energy $E$ are plotted as functions of $(x^+, x^-)$. Note that the energy
is positive and tends to zero at infinity.
%In perturbations, the sign of the initial Bondi energy $E_+(0)=E(\infty,0)$
%will determine the fate of the wormhole. See Table \ref{table1}.
}
\label{fig3_config}
\end{figure}
%<<<<<<<<<<<<<<<<<<<<<<<<<<<<<<<<<<<<<<<     \ref{fig3_config}


%>>>>>>>>>>>>>>>>>>>>>>>>>>>>>>>>>>>>>>> Fig.\ref{fig4_pulseGexp}
\begin{figure}[h]
\setlength{\unitlength}{1mm}
\begin{picture}(85,70)
\put(10,0){\epsfxsize=60mm \epsffile{fig4a.eps} }
\end{picture}
\caption[quartic]{ Horizon locations, $\vartheta_\pm=0$, for perturbed
wormhole. Fig.(a) is the case we supplement the ghost field, $c_a = 0.1$, and
(b1) and (b2) are where we reduce the field, $c_a = -0.1$ and $-0.01$. Dashed
lines and solid lines are $\vartheta_+=0$ and $\vartheta_-=0$ respectively. In
all cases, the pulse hits the wormhole throat at $(x^+, x^-)=(3,3)$. A
$45^\circ$ counterclockwise rotation of the figure corresponds to a partial
Penrose diagram. } \label{fig4_pulseGexp}
\end{figure}
%<<<<<<<<<<<<<<<<<<<<<<<<<<<<<<<<<<<<<<<     \ref{fig4_pulseGexp}

\setcounter{figure}{3}
%>>>>>>>>>>>>>>>>>>>>>>>>>>>>>>>>>>>>>>> Fig.\ref{fig4_pulseGexp}
\begin{figure}[h]
\setlength{\unitlength}{1mm}
\begin{picture}(85,140)
\put(10,70) {\epsfxsize=60mm \epsffile{fig4b1.eps} }
\put(10,0)  {\epsfxsize=60mm \epsffile{fig4b2.eps} }
\end{picture}
\caption[quartic]{  (continued) }
%\label{fig4_pulseGexp}
\end{figure}
%<<<<<<<<<<<<<<<<<<<<<<<<<<<<<<<<<<<<<<<     \ref{fig4_pulseGexp}


%>>>>>>>>>>>>>>>>>>>>>>>>>>>>>>>>>>>>>>> Fig.\ref{fig5_sketch}
\begin{figure}[h]
\setlength{\unitlength}{1mm}
\begin{picture}(80,60)
\put(5,0){\epsfxsize=65mm \epsffile{fig5.eps} }
\end{picture}
\caption[quartic]{Partial Penrose diagram of the evolved space-time. }
\label{fig5_sketch}
\end{figure}
%<<<<<<<<<<<<<<<<<<<<<<<<<<<<<<<<<<<<<<<     \ref{fig5_sketch}


%>>>>>>>>>>>>>>>>>>>>>>>>>>>>>>>>>>>>>>> Fig.\ref{fig6_throat}
\begin{figure}[t]
\setlength{\unitlength}{1mm}
\begin{picture}(85,60)
\put(5,0){\epsfxsize=70mm \epsffile{fig6.eps} }
%\put(5,0){\epsfxsize=75mm \epsffile{fig7b.eps} }
\end{picture}
\caption[quartic]{
%(a)
Areal radius $r$ of the ``throat'' $x^+=x^-$, plotted as a function of proper
time. Additional negative energy causes inflationary expansion, while reduced
negative energy causes collapse to a black hole and central singularity.
  } \label{fig6_throat}
\end{figure}
%<<<<<<<<<<<<<<<<<<<<<<<<<<<<<<<<<<<<<<<     \ref{fig6_throat}


%>>>>>>>>>>>>>>>>>>>>>>>>>>>>>>>>>>>>>>> Fig.\ref{fig7_localE}
\begin{figure}[t]
\setlength{\unitlength}{1mm}
\begin{picture}(85,160)
\put(5,110){\epsfxsize=65mm \epsffile{fig7a.eps} }
\put(5,55) {\epsfxsize=65mm \epsffile{fig7b1.eps} }
\put(5,0)  {\epsfxsize=65mm \epsffile{fig7b2.eps} }
\end{picture}
\caption[quartic]{ Energy $E(x^+,x^-)$ as a function of $x^-$, for  $x^+ =
12,16,20$. Here $c_a$ is (a) $0.05$, (b1) $-0.1$ and (b2) $-0.01$. The energy
for different $x^+$ coincides at the final horizon location $x^-_H$, indicating
that the horizon quickly attains constant mass $M=E(\infty,x^-_H)$. This is the
final mass of the black hole. The values are shown in Table~\ref{table1}.
  } \label{fig7_localE}
\end{figure}
%<<<<<<<<<<<<<<<<<<<<<<<<<<<<<<<<<<<<<<<     \ref{fig7_localE}

%>>>>>>>>>>>>>>>>>>>>>>>>>>>>>>>>>>>>>>> Fig.\ref{fig8_BHmass}
\begin{figure}[t]
\setlength{\unitlength}{1mm}
\begin{picture}(85,120)
\put(0,60){\epsfxsize=70mm \epsffile{fig8a.eps} }
\put(0,0) {\epsfxsize=70mm \epsffile{fig8b.eps} }
\end{picture}
\caption[quartic]{ Relation between the initial perturbation and the final mass
of the black hole. (a) The final horizon ($\vartheta_+=0$) coordinate
$x^-_H-3$, since we fixed $c_c=3$, versus initial energy of the perturbation,
$E_0$. We plotted the results of the runs of $c_a=10^{-1}, \cdots, 10^{-4}$
with $c_b=3, 6,$ and 9. They lie close to one line. (b) The final black hole
mass $M$ for the same examples. We see that $M$ appears to reach a non-zero
minimum for small perturbations.
  } \label{fig8_BHmass}
\end{figure}
%<<<<<<<<<<<<<<<<<<<<<<<<<<<<<<<<<<<<<<<     \ref{fig8_BHmass}



%>>>>>>>>>>>>>>>>>>>>>>>>>>>>>>>>>>>>>>> Fig.\ref{fig9_pulseNexp}
\begin{figure}[bht]
\setlength{\unitlength}{1mm}
\begin{picture}(85,70)
\put(5,0) {\epsfxsize=60mm \epsffile{fig9.eps} }
\end{picture}
\caption[quartic]{
Evolution of a wormhole perturbed by a normal scalar field.
Horizon locations:  dashed lines and solid lines are $\vartheta_+=0$
and $\vartheta_-=0$ respectively.
}
\label{fig9_pulseNexp}
\end{figure}
%<<<<<<<<<<<<<<<<<<<<<<<<<<<<<<<<<<<<<<<     \ref{fig9_pulseNexp}

%>>>>>>>>>>>>>>>>>>>>>>>>>>>>>>>>>>>>>>> Fig.\ref{fig10_BHmass}
\begin{figure}[t]
\setlength{\unitlength}{1mm}
\begin{picture}(85,120)
\put(0,60){\epsfxsize=70mm \epsffile{fig10a.eps} }
\put(0,0) {\epsfxsize=70mm \epsffile{fig10b.eps} }
\end{picture}
\caption[quartic]{ The same plots as in Fig.\ref{fig8_BHmass} for conventional
field pulses. (a) The final horizon ($\vartheta_+=0$) coordinate $x^-_H-3$
versus initial energy of the perturbation, $E_0$. We plotted the results of the
runs of $\tilde{c}_a=0.5, \cdots, 10^{-2}$ with $\tilde{c}_b=3, 6,$ and 9. They
lie close to one line, different to that of Fig.\ref{fig8_BHmass}(a), but with
the same gradient. (b) The final black hole mass $M$ for the same examples.
  } \label{fig10_BHmass}
\end{figure}
%<<<<<<<<<<<<<<<<<<<<<<<<<<<<<<<<<<<<<<<     \ref{fig10_BHmass}



%>>>>>>>>>>>>>>>>>>>>>>>>>>>>>>>>>>>>>>> Fig.\ref{fig11_maint}
\begin{figure}[bht]
\setlength{\unitlength}{1mm}
\begin{picture}(85,70)
\put(5,0) {\epsfxsize=60mm \epsffile{fig11.eps} }
\end{picture}
\caption[quartic]{ Temporary wormhole maintenance.  After a normal scalar pulse
representing a traveller, we beamed in an additional ghost pulse to extend the
life of the wormhole. Horizon locations $\vartheta_+=0$ are plotted for three
cases: (A) no maintenance, which results in a black hole; (B) with a
maintenance pulse which results in an inflationary expansion; (C) with a more
finely tuned maintenance pulse, which keeps the static structure up to the end
of the range. } \label{fig11_maint}
\end{figure}
%<<<<<<<<<<<<<<<<<<<<<<<<<<<<<<<<<<<<<<<     \ref{fig11_maint}
